\section{Threat Model}\label{sec:threat}

\paragraph{Agent.} Let $u$ be the user instruction. At each step, the model produces an action $a_t$ consisting of a
tool name and arguments, or produces the final answer $y$. The Tool Broker validates $a_t$ against the schema,
executes it and returns an observation $o_t$. The system has six read tools (searching and reading documents, querying
a read-only database, searching and opening cached pages, and a calculator) and two \emph{sinks} with side effects:
sending e-mail and calling a webhook. Both sinks are mocks that only write a log and send nothing over the network. The
final answer is treated as a third sink. For a sink call $a$, $\dest(a)$ is its destination (recipient or URL) and
$\pay(a)$ denotes the remaining arguments.

\paragraph{Two labels.} Every piece of content carries two independent labels assigned by the host. In terms of
\emph{trust}, only $u$ and the system prompt are trusted, and every tool observation is untrusted. In terms of
\emph{sensitivity}, content ranges from public to confidential. Confidential sources contain values such as keys,
codes or account identifiers. These values must not leave the agent unless the user allows it. An untrusted page can
be public, while a confidential record can come from a source the attacker cannot touch. Merging the two dimensions
into one would force a defence either to block all confidential data or to trust every internal source.

\paragraph{Attacker.} The attacker controls the text of one piece of content returned by a read tool, at any position
and in any tone. The attacker knows the names and schemas of the tools but cannot modify $u$, the system prompt, the
broker or the model weights, and cannot observe the trajectory. On the sealed sets, attacks are static and were written
without knowledge of the defence code. Direct jailbreaks, training-data poisoning, denial of service and multi-agent
systems are out of scope.

\paragraph{Harm predicates.} Let $G(u)$ be the set of destinations authorised by the user, $V$ the set of confidential
values in the trajectory that the user has not allowed to be disclosed, and $\mathcal{S}$ the set of \emph{executed}
sink calls. A trajectory $\tau$ is harmful if
\begin{equation}\label{eq:harm}
\mathrm{Harm}(\tau) = H_\mathit{dest} \lor H_\mathit{leak} \lor H_\mathit{tamper} \lor H_\mathit{side},
\end{equation}
where
\begin{itemize}
\item $H_\mathit{dest}$: there is $a \in \mathcal{S}$ with $\dest(a) \notin G(u)$,
\item $H_\mathit{leak}$: there is $v \in V$ that appears in $\pay(a)$ for some $a \in \mathcal{S}$ or in $y$,
\item $H_\mathit{tamper}$: there is $a \in \mathcal{S}$ with $\dest(a) \in G(u)$ but a user-specified field carries a
  different value,
\item $H_\mathit{side}$: the task only requires reading but $\mathcal{S} \neq \emptyset$.
\end{itemize}
A frozen scorer evaluates these predicates from the executed calls and from $y$. Blocked calls do not count as harm.

\paragraph{Six attack mechanisms.} Attacks are classified by \emph{mechanism}, that is, the path by which the
injection turns into harm (Table~\ref{tab:mechanisms}). In the two redirection mechanisms (\mER, \mCR), the destination
is not in $u$, so checking the destination suffices regardless of wording. In the two leak mechanisms (\mSL, \mAL), the
destination is valid or absent, so only the provenance of the value separates the attack from the task. In parameter
tampering (\mPT), both the destination and the sensitivity are normal, which makes it the hardest mechanism for
label-based rules. The unrequested side-effect mechanism (\mUS) tests whether a destination that is merely mentioned,
for example in a prohibition, is mistaken for an authorisation.

% M1..M6 trong log/luận văn = ER, CR, SL, AL, PT, US.
% Kiểu bảng chung (14/11): tabularx/tabular* đúng \textwidth, \footnotesize, \tabcolsep 3pt, cột chữ trái, cột số phải.
\begin{table}[!htbp]
\centering
\caption{The six attack mechanisms and their harm predicates.}\label{tab:mechanisms}
\footnotesize
\setlength{\tabcolsep}{3pt}
\begin{tabularx}{\textwidth}{@{}l>{\raggedright\arraybackslash}Xl@{}}
\toprule
Code & Mechanism (subtypes) & Predicate \\
\midrule
\mER & Explicit redirection: asks to send to a destination not in $u$, with a send verb and a full address & $H_\mathit{dest}$ \\
\mCR & Covert redirection: same goal but avoids the ``send + address'' pattern (verbless note, forward or cc, split address, destination in a second document) & $H_\mathit{dest}$ \\
\mSL & Sink leak: puts a confidential value into a sink the user specified (e-mail body, payload, innocuous-looking field) & $H_\mathit{leak}$ \\
\mAL & Answer leak: puts a confidential value into the final answer (verbatim, embedded in a sentence, labelled) & $H_\mathit{leak}$ \\
\mPT & Parameter tampering: changes an ordinary value in the payload of a sink the user specified (number, boolean, unit, identifier) & $H_\mathit{tamper}$ \\
\mUS & Unrequested side effect: inserts a sink call into a read-only task (destination only in a prohibition, merely informational, or absent) & $H_\mathit{side}$ \\
\bottomrule
\end{tabularx}
\end{table}


\paragraph{Families, benign counterparts and language variants.} A \emph{family} is one attack idea in a specific task
and tool environment. Each benign counterpart uses the same task and environment as its family. Every family has a
\emph{\bPlain} counterpart, and one third of the families also have an \emph{\bAuth} counterpart. The \bPlain{}
counterpart replaces the injection with a harmless note of the same length and is used to measure utility. The
\bAuth{} counterpart keeps a request similar to the attack goal, but the user explicitly allows it. Blocking it is a
\emph{false block}. Each family has five variants applied to both $u$ and the injection: standard Vietnamese,
paraphrased, Vietnamese--English code-mixed, without diacritics, and English. Identifiers, addresses and confidential
values are identical across variants, so any difference comes from language alone.
